\begin{abstract}
We prove that the rescaled counting measure of a critical square-lattice
Fortuin--Kasteleyn Dobrushin interface in the unit square converges to the
Minkowski-content measure of its Schramm--Loewner limit, for every fixed
cluster weight $1\le q<4$. A single deterministic constant times the
predicted power of the mesh gives the normalization. Convergence is joint
with the ordered curve and includes the total mass, giving the scaling
limit of the interface's total number of steps.
\end{abstract}

\section{Introduction}\label{sec:introduction}

A scaling limit of an interface describes its shape, but does not by
itself describe how many lattice steps it takes to traverse that shape.
The latter question requires a measure on the limiting curve. For
Schramm--Loewner evolution (SLE), the appropriate measure is its
dimensional Minkowski content. The problem here is to identify that
measure as the limit of the unweighted number of steps of a critical
Fortuin--Kasteleyn (FK) interface on the square lattice.

The distinction between shape and occupation is quantitative. A
visitation exponent predicts the appropriate power of the mesh, but
does not rule out a slowly varying or nonconvergent factor. We prove
convergence with a single constant multiplying the power law. We also
control the mass near the boundary, so the result concerns the whole
interface, including its total number of steps.

\subsection{The lattice model and the main result}

Fix $q\in[1,4)$ and put
\begin{equation}\label{eq:parameters}
 p_q=\frac{\sqrt q}{1+\sqrt q},\qquad
 \kappa=\frac{4\pi}{\arccos(-\sqrt q/2)},\qquad
 d=1+\frac{\kappa}{8},\qquad x=2-d.
\end{equation}
Thus $4<\kappa\le6$ and $1/4\le x<1/2$.
Let $D=(0,1)^2$, with marked corners $a=(0,0)$ and $b=(1,1)$.
The counterclockwise boundary arc from $a$ to $b$, consisting of the
bottom and right sides, is denoted by $F$. The complementary arc is $W$.

For $n\ge1$, the graph has vertex set
$n^{-1}\Z^2\cap\overline D$ and all nearest-neighbor edges in
$\overline D$. Identify all vertices on $W$, including $a,b$, in a
single block; all other vertices are distinct. Denote this partition by
$\xi$. Remove the boundary edges on $W$ from the sampled set of edges,
and call the remaining set $E_n'$. The probability of a set of open
edges $A\subset E_n'$ is
\[
 \Pp_n(A)=\frac{1}{Z_n}
 p_q^{|A|}(1-p_q)^{|E_n'|-|A|}q^{k_\xi(A)},
\]
where $k_\xi(A)$ counts connected components after the prescribed
identification, including isolated vertices. The removed edges on $W$
are declared open for drawing the interface. Edges on $F$ remain random.

We use the Dobrushin medial strand from $a$ to $b$. Medial vertices are
midpoints of primal edges; medial edges join midpoints of consecutive
sides of a lattice square. At an open primal edge, pair the four medial
half-edges in the two turns around the adjacent dual faces. At a closed
primal edge, use the turns around the primal endpoints. The boundary
turns separate the primal wired arc from the complementary dual wired
arc. Two distinguished terminal half-edges occur at the boundary changes.
The pairing determines the strand and its order, including the order at
a medial vertex visited twice.

Let $e_{n,1},\ldots,e_{n,N_n}$ be the full medial-edge traversals of
this strand, including boundary traversals but excluding the two
terminal half-edges. Let $z_{n,j}$ be the midpoint of $e_{n,j}$.
Define
\[
 \Pi(u,v)=\bigl(\max(0,\min(1,u)),\max(0,\min(1,v))\bigr).
\]
Project the drawn strand by $\Pi$ and join its terminal points to $a,b$
by straight segments. The resulting oriented curve is $\eta_n$.
The terminal connectors lie within $O(n^{-1})$ of the marked corners
and carry no counting mass. For $c>0$, set
\begin{equation}\label{eq:counting-measure}
 \nu_{n,c}=c n^{-d}\sum_{j=1}^{N_n}\delta_{\Pi(z_{n,j})}.
\end{equation}
This convention counts traversals, not Euclidean length. In fact, a full
medial edge cannot recur on this strand: after resolving the pairings,
the strand is a component of a graph of degree two, with its two terminal
ends. A medial vertex can nevertheless be visited twice. This distinction
will matter when we replace counting by local hit indicators.

For continuous oriented curves $\gamma_1,\gamma_2:[0,1]\to\overline D$,
use the distance
\begin{equation}\label{eq:curve-metric}
 d_{\mathrm{curv}}(\gamma_1,\gamma_2)
 =\inf_{\phi_1,\phi_2}\sup_{0\le t\le1}
 |\gamma_1(\phi_1(t))-\gamma_2(\phi_2(t))|,
\end{equation}
where the infimum is over increasing homeomorphisms of $[0,1]$.
We identify curves at distance zero. This topology retains traversal
order, unlike a topology on trace sets alone.

Let $\eta$ be chordal $\mathrm{SLE}_\kappa$ in $(D;a,b)$, oriented
from $a$ to $b$. Its $d$-dimensional Minkowski-content measure is
normalized by
\begin{equation}\label{eq:content-measure}
 \mu_{\eta,r}
 =r^{-x}\Pi_*\bigl[
 \1\{\dist(z,\tr\eta)<r\}\,\mathrm dA(z)\bigr],
 \qquad \mu_{\eta,r}\longrightarrow\mu_\eta
 \quad\text{in probability as }r\downarrow0.
\end{equation}
Here the area measure is on the whole plane, not only on $D$.
The convergence is weak convergence of finite measures on
$\overline D$. The existence of this whole-trace measure, including
the stated projection convention, is justified in
Lemma~\ref{lem:whole-content} from the SLE content and Green-function
theorems.

\begin{theorem}\label{thm:main}
For each fixed $q\in[1,4)$ there is a deterministic constant
$c(q)\in(0,\infty)$ such that, along all positive integers $n$,
\[
 \bigl([\eta_n],\nu_{n,c(q)}\bigr)
 \ \Longrightarrow\
 \bigl([\eta],\mu_\eta\bigr).
\]
The first coordinate has the ordered curve topology
\eqref{eq:curve-metric}, and the second has the weak topology of finite
measures on $\overline D$. The two limiting coordinates belong to the
same SLE sample. In particular,
\[
 c(q)n^{-d}N_n\ \Longrightarrow\ \mu_\eta(\overline D)
\]
jointly with the curve and all continuous spatial tests of its measure.
\end{theorem}

The value of $c(q)$ depends on the stated discrete and continuum units;
its numerical evaluation is not needed. Every estimate below is for a
fixed $q$. We make no claim of uniformity as $q\uparrow4$.

\subsection{Context and inputs}

Schramm introduced SLE to describe conformally invariant limits of
planar lattice interfaces \cite{Schramm}. Beffara proved that the SLE
trace has dimension $1+\kappa/8$ for $\kappa\le8$ \cite{Beffara}.
Dimension predicts the exponent in \eqref{eq:counting-measure}; it
does not identify a counting-measure limit. Lawler and Sheffield
introduced a Doob--Meyer construction of natural parametrization
\cite{LawlerSheffield}; Lawler and Zhou established its existence
throughout $0<\kappa<8$ \cite{LawlerZhou}. Lawler and Rezaei proved
that the resulting natural measure agrees, up to normalization, with
$d$-dimensional Minkowski content \cite{LawlerRezaeiContent}. We use
their interior content theorem and Euclidean Green-function estimates
\cite{LawlerRezaeiBasic}, and supply the whole-square boundary argument
in Section~\ref{sec:occupation}.

There are two particularly relevant precedents for passing from a
lattice curve to its natural measure. For critical triangular-lattice
site percolation, Garban, Pete, and Schramm constructed an interface
measure measurable from the limiting SLE$_6$ curve, with conformal
covariance exponent $7/4$ \cite{GPS}. Their normalization at mesh
$\eta$ is $\eta^2/\alpha_2^\eta(\eta,1)$, where
$\alpha_2^\eta(\eta,1)$ is the alternating two-arm probability.
Holden, Li, and Sun identified the resulting measure with Minkowski
content and proved convergence in natural parametrization \cite{HLS}.
The later sharp alternating-arm asymptotics of Du, Gao, Li, and Zhuang
\cite{DGLZ} permit that normalization to be replaced by a constant
times $\eta^{7/4}$. These are triangular-site percolation results,
not occupation results for square-lattice FK interfaces.

For square-lattice loop-erased random walk, Lawler and Viklund proved
convergence to SLE$_2$ when each step receives a deterministic constant
times the $5/4$ power of the mesh \cite{LawlerViklund}. Their result
uses model-specific Green-function estimates and is stated for bounded
simply connected domains with analytic boundary. In the FK setting,
the Ising case already has ordered-trace convergence to SLE$_{16/3}$
\cite{IsingInterfaces}. These precedents distinguish the geometric
trace limit from the microscopic estimates required for natural
measure convergence.

For the present FK model we take chordal ordered-trace convergence
from the companion manuscript \cite[Theorem \texttt{chordal:trace}]{FKCompanion}.
This is the trace input specified in Input~\ref{inp:trace}. We separately
identify the additional companion results used in the argument:
buffered FK comparison and passage estimates, the whole-plane loop
limit, and the square-row representation. They are cited at their
stated scopes in Section~\ref{sec:inputs}; none supplies the occupation
asymptotic proved here. The public crossing and arm-separation inputs
are due to Duminil-Copin, Manolescu, and Tassion \cite{DMT}, building
on the FK-Ising separation arguments of Chelkak, Duminil-Copin, and
Hongler \cite{CDCH}. The electric normalization is fixed by the
six-vertex height covariance \cite{DKLM,GFFCompanion}.

One parameter issue requires care. The angular-current proofs in the
GFF companion are presented in a parameter interval different from
the FK interval. In Section~\ref{sec:spectral} we apply their algebraic
arguments to the FK companion's row representation and verify the
required hypotheses there. We do not import an angular-spectrum
theorem outside its stated range.

\subsection{How the normalization is determined}

There are two distinct obstacles. First, one must obtain information
more precise than an arm exponent. Second, one must compare that
information with visits to a single medial edge. The proof separates
these tasks.

We orient loops and insert phases depending on their winding around
marked lattice points. Summing the orientations gives scalar loop
weights; for small charges these weights are positive. Regeneration
across surrounding open circuits gives quantitative loss of dependence
on the exterior and zero-free complex neighborhoods for the associated
partition functions. These estimates are developed in
Section~\ref{sec:radial}.

The height potential associated with small electric charges is nearly
harmonic, with a summable power error. The proof uses two directions
of row slicing: the untwisted current controls the potential away from
the line of insertions, and twisted rows control its differences near
that line. Step-two lattice potential theory then gives a Coulomb
correlation formula with a power remainder. Complex continuation
extends it to charges that impose arm events, and a reflection symmetry
removes the remaining linear ambiguity in the local normalization.
Sections~\ref{sec:spectral} and~\ref{sec:electric} contain this part of
the argument.

To reach a single lattice edge, we condition two opposite-color arms
to reach a large radius. Section~\ref{sec:coupling} constructs a
coupling that makes the exterior configurations identical with power
accuracy. The successful separating band has exactly two medial
strands crossing it, so its exterior continuation is independent of
the microscopic starting configuration. This last property is stronger
than separation of two chosen arms; it is what makes the coupling
useful for counting a distinguished interface.

Section~\ref{sec:calibration} compares two electric charges at a
mesoscopic separation with charges at neighboring lattice points.
The coupling applies even though the remaining electric factor is
signed. A nonzero mesoscopic amplitude controls the relative error,
and a summable comparison between scales $n$ and $N\in[n,2n]$ gives a
genuine microscopic amplitude. A second comparison, with a disk probe,
produces deterministic constants for replacing lattice visits by
neighborhood area.

Finally, Section~\ref{sec:occupation} performs this replacement in
$L^2$. Ordered-trace convergence identifies the disk probes, while
SLE Minkowski content identifies their small-radius limit. Integrable
boundary-strip estimates complete the passage to measures on the
closed square.

The potentially reusable ingredients are the complex stability of
radially tilted loop partition functions, the two-direction estimate
for the electric height potential, and the separating-band coupling
used to convert mesoscopic correlation amplitudes into microscopic
normalization. Their roles are quantitative: they remove errors that
an exponent-only argument cannot detect.
